\chapter{Introduction}\label{sec:introduction}

\section{Cooperative Relay Communication}\label{sec:cooperative-relay-communication}

Cooperative relay communication is a fundamental technique in modern wireless networks that extends coverage, improves reliability, and increases throughput by employing intermediate nodes between a source and destination. The theoretical study of relay channels dates to van der Meulen's formulation in 1971 and the seminal capacity bounds of Cover and El Gamal \cite{cover1979capacity}, who established inner and outer bounds for the general relay channel that remain the tightest known results for many configurations. In a two-hop relay network, a source transmits a signal to a relay, which processes it and retransmits it to the destination. This architecture is central to standards such as LTE-Advanced and 5G NR, where relay nodes bridge coverage gaps and enhance cell-edge performance \cite{laneman2004cooperative}, \cite{nosratinia2004cooperative}.

\subsection{Information-Theoretic Foundations}\label{sec:information-theoretic-foundations}

The three-terminal relay channel consists of a source \(S\), a relay \(R\), and a destination \(D\). The capacity of this channel depends critically on the relay processing strategy \(f(\cdot)\) and the channel statistics. Cover and El Gamal \cite{cover1979capacity} established two fundamental capacity bounds:

\textbf{Cut-set upper bound.} The capacity of any relay channel is bounded by:

\begin{equation}\protect\phantomsection\label{eq:relay-capacity-bound}{C \leq \max_{p(x, x_R)} \min\left\{I(X, X_R; Y_D), \; I(X; Y_R, Y_D | X_R)\right\}
}\end{equation}\par{\small\textit{Source: \cite{cover1979capacity, elgamal2011network}}}

where the first term represents the maximum rate at which the destination can receive information from both the source and relay jointly (the broadcast cut), and the second term represents the maximum rate at which the source can communicate to both the relay and destination (the multiple-access cut). The capacity is limited by the weaker of these two cuts, establishing a fundamental bottleneck principle for relay communication.

\textbf{DF achievability bound.} Under decode-and-forward, the relay fully decodes the source message and cooperates with the source to transmit a common message to the destination:

\begin{equation}\protect\phantomsection\label{eq:df-capacity}{C_{\text{DF}} = \max_{p(x, x_R)} \min\left\{I(X; Y_R | X_R), \; I(X, X_R; Y_D)\right\}
}\end{equation} \par{\small\textit{Source: \cite{cover1979capacity}}}

This rate is achievable using block Markov encoding and backward decoding. The first term is the relay's decoding constraint (it must fully decode the source), and the second is the destination's decoding rate. DF achieves the cut-set bound when the source-relay link is stronger than the relay-destination link, making it capacity-optimal for near-source relays.

For the degraded relay channel --- where the relay's observation is a degraded version of the destination's, or vice versa --- Cover and El Gamal showed that the DF bound coincides with the cut-set bound, establishing the exact capacity. In the Gaussian case with equal per-hop SNR \(\gamma\), the two-hop DF capacity is:

\begin{equation}\protect\phantomsection\label{eq:df-gaussian-capacity}{C_{\text{DF}}^{\text{Gaussian}} = \frac{1}{2}\log_2\left(1 + \gamma\right)
}\end{equation} \par{\small\textit{Source: \cite{laneman2004cooperative, cover1979capacity}}}

which equals the single-hop capacity --- i.e., the relay incurs no rate penalty when the source-relay link is sufficiently strong.

\subsection{Two-Hop Relay Model}\label{sec:two-hop-relay-model}

In the half-duplex two-hop model studied in this thesis, the relay cannot transmit and receive simultaneously, and there is no direct source-destination link. The communication proceeds in two time slots:

\textbf{Slot 1 (Source → Relay):} \begin{equation}\protect\phantomsection\label{eq:awgn-hop1}{y_R = x + n_1, \quad n_1 \sim \mathcal{N}(0, \sigma^2)
}\end{equation} \par{\small\textit{Source: \cite{laneman2004cooperative}}}

\textbf{Slot 2 (Relay → Destination):} \begin{equation}\protect\phantomsection\label{eq:awgn-hop2}{y_D = x_R + n_2, \quad n_2 \sim \mathcal{N}(0, \sigma^2)
}\end{equation} \par{\small\textit{Source: \cite{laneman2004cooperative}}}

where \(x \in \{-1, +1\}\) is the transmitted BPSK symbol, \(y_R\) is the signal received at the relay, \(x_R = f(y_R)\) is the relay's output after processing, and \(y_D\) is the signal received at the destination. The noise terms are independent and identically distributed with variance \(\sigma^2 = P_s / \text{SNR}\). The half-duplex constraint introduces a spectral efficiency penalty of factor 2 (since two time slots are needed per symbol), but this is a common assumption in practical relay standards and simplifies the analysis while remaining representative of many practical relay-processing scenarios.

The choice of relay processing function \(f(\cdot)\) fundamentally determines system performance and is the central object of study in this thesis. Classical approaches include amplify-and-forward (AF), which simply scales the received signal, and decode-and-forward (DF), which regenerates the signal through demodulation and re-modulation. Each has well-understood performance characteristics: AF is simple but propagates noise, while DF eliminates first-hop noise but introduces error propagation when decoding fails \cite{cover1979capacity}.

\subsection{Cooperative Diversity and Practical Relevance}\label{sec:cooperative-diversity-and-practical-relevance}

Laneman, Tse, and Wornell \cite{laneman2004cooperative} showed that cooperative relaying achieves spatial diversity without requiring multiple antennas at any single node. Specifically, a system with \(L\) cooperating single-antenna relays can achieve a diversity order of \(L + 1\), meaning the outage probability decays as \(\text{SNR}^{-(L+1)}\). For a single relay (\(L=1\)), this yields second-order diversity --- a significant improvement over the first-order diversity of point-to-point communication over fading channels.

In practical deployments, relay nodes serve several roles: (i) \textbf{range extension} for cell-edge users where the direct link is too weak, (ii) \textbf{coverage filling} in shadowed areas behind buildings or terrain, (iii) \textbf{capacity enhancement} through spatial reuse, and (iv) \textbf{energy efficiency} by reducing transmission power requirements through shorter per-hop distances. The 3GPP LTE-Advanced standard (Release 10+) defines Type I (non-transparent) and Type II (transparent) relay nodes, while 5G NR introduces integrated access and backhaul (IAB) nodes that extend this concept to millimeter-wave and sub-THz bands. The AI-based relay processing studied in this thesis is applicable to any of these deployment scenarios, as the neural network operates at the baseband processing level independently of the RF front-end and protocol layer.

The fundamental question that motivates this thesis is: can a learned relay function \(f_\theta(\cdot)\) outperform the classical relay functions (AF scaling, DF regeneration) by exploiting patterns in the received signal that analytical methods do not capture? This question is particularly relevant at low SNR, where both AF (noise amplification) and DF (decoding errors) have well-known limitations, and where the non-linear decision boundary learned by a neural network may provide a substantive advantage.

\section{Classical Relay Strategies}\label{sec:classical-relay-strategies}

\textbf{Amplify-and-Forward (AF).} The AF relay amplifies the received signal by a gain factor \(G\) that normalizes the output power:

\begin{equation}\protect\phantomsection\label{eq:af-gain}{G = \sqrt{\frac{P_{\text{target}}}{\mathbb{E}[|y_R|^2]}}
}\end{equation} \par{\small\textit{Source: \cite{laneman2004cooperative}}}

The end-to-end SNR for AF relay is:

\begin{equation}\protect\phantomsection\label{eq:af-snr}{\text{SNR}_{\text{eff}}^{\text{AF}} = \frac{\text{SNR}_1 \cdot \text{SNR}_2}{\text{SNR}_1 + \text{SNR}_2 + 1}
}\end{equation} \par{\small\textit{Source: \cite{laneman2004cooperative}}}

This expression reveals a fundamental limitation: even when one hop has high SNR, the effective SNR is bottlenecked by the weaker hop. Moreover, the relay amplifies both signal and noise from the first hop, leading to noise accumulation at the destination.

\textbf{Decode-and-Forward (DF).} The DF relay demodulates the received signal, recovers the transmitted bits, and re-modulates clean symbols:

\begin{equation}\protect\phantomsection\label{eq:df-demod-remod}{\hat{b}_R = \text{demod}(y_R), \quad x_R = \text{mod}(\hat{b}_R)
}\end{equation} \par{\small\textit{Source: \cite{laneman2004cooperative}}}

The end-to-end BER for DF is:

\begin{equation}\protect\phantomsection\label{eq:df-ber-hops}{P_e^{\text{DF}} = P_{e,1} + (1 - P_{e,1}) \cdot P_{e,2}
}\end{equation} \par{\small\textit{Source: \cite{laneman2004cooperative, proakis2008digital}}}

where \(P_{e,1}\) and \(P_{e,2}\) are the BER of the first and second hops, respectively. For BPSK modulation over AWGN (real-valued noise with variance \(1/\text{SNR}\)), each hop's BER is \(P_e = Q\left(\sqrt{\text{SNR}}\right)\). DF provides clean regeneration at high SNR but suffers from error propagation when the first-hop BER is non-negligible.

\textbf{Other Classical Relay Techniques.} Beyond AF and DF, several additional cooperative relay strategies have been proposed in the literature. \textbf{Compress-and-Forward (CF)} has the relay quantize and compress the received signal using Wyner-Ziv source coding, then forward the compressed representation to the destination, which exploits its own direct-link observation as side information to decode \cite{cover1979capacity}, \cite{elgamal2011network}. \textbf{Compute-and-Forward (CoF)} leverages lattice codes to allow relays to decode \emph{linear combinations} of transmitted messages rather than individual codewords; the destination then solves a system of linear equations to recover the original messages \cite{nazer2011compute}. \textbf{Estimate-and-Forward (EF)} computes an MMSE estimate of the transmitted signal at the relay and forwards the soft estimate, avoiding the hard-decision errors of DF while providing a more structured output than AF. \textbf{Selective and incremental relaying} are hybrid protocols: in selective relaying the relay forwards only when it can decode correctly, while in incremental relaying the relay transmits only when the destination signals (via a feedback link) that the direct transmission has failed \cite{laneman2004cooperative}. \textbf{Two-way relaying with network coding} allows both endpoints to transmit simultaneously to the relay, which combines the received signals (e.g., via XOR or lattice coding) and broadcasts the result back, effectively doubling spectral efficiency \cite{rankov2007spectral}. \textbf{Filter-and-Forward (FF)} applies a linear filter (e.g., matched filter or Wiener filter) at the relay before forwarding, providing a structured intermediate between AF's simple scaling and DF's full regeneration.

AF and DF represent the two canonical extremes of the relay processing spectrum: AF performs minimal processing (linear scaling) while DF performs maximal classical processing (full regeneration). The other techniques --- CF, EF, CoF --- occupy intermediate points along this spectrum. In this thesis, the AI-based relay strategies are designed to \emph{learn} the optimal relay processing function from data, potentially discovering strategies that subsume or outperform these fixed classical schemes. The focus on AF and DF as classical baselines is therefore deliberate: they bracket the classical design space and provide clear lower and upper bounds against which the learned strategies can be evaluated. These additional relay strategies are reviewed for completeness; they are not simulated in this thesis because the studied system assumes a two-hop link without a direct source--destination path.

\section{MIMO Systems and Equalization}\label{sec:mimo-systems-and-equalization}

Multiple-input multiple-output (MIMO) systems employ multiple antennas at both transmitter and receiver to exploit spatial multiplexing and diversity gains \cite{tse2005fundamentals}. The theoretical foundation was established by Telatar \cite{telatar1999capacity} and Foschini \cite{foschini1996layered}, who independently showed that the ergodic capacity of an \(N_t \times N_r\) MIMO channel with independent Rayleigh fading scales as:

\begin{equation}\protect\phantomsection\label{eq:mimo-capacity}{C = \mathbb{E}\left[\log_2 \det\left(\mathbf{I}_{N_r} + \frac{\text{SNR}}{N_t}\mathbf{H}\mathbf{H}^H\right)\right]
}\end{equation} \par{\small\textit{Source: \cite{telatar1999capacity, foschini1996layered}}}

For the \(2 \times 2\) system used in this thesis, this yields approximately \(C \approx 2\log_2(1 + \text{SNR}/2)\) at high SNR --- a doubling of the SISO capacity, achieved by transmitting independent data streams on each antenna.

\subsection{System Model}\label{sec:system-model}

In a \(2 \times 2\) MIMO system, the received signal is:

\begin{equation}\protect\phantomsection\label{eq:mimo-received}{\mathbf{y} = \mathbf{H}\mathbf{x} + \mathbf{n}
}\end{equation} \par{\small\textit{Source: \cite{tse2005fundamentals}}}

where \(\mathbf{H} \in \mathbb{C}^{2 \times 2}\) is the channel matrix with \(H_{ij} \sim \mathcal{CN}(0, 1)\) (independent Rayleigh fading per link), \(\mathbf{x} \in \mathbb{C}^{2}\) is the transmitted symbol vector with \(\mathbb{E}[\mathbf{x}\mathbf{x}^H] = (P/N_t)\mathbf{I}\), and \(\mathbf{n} \sim \mathcal{CN}(\mathbf{0}, \sigma^2\mathbf{I})\) is noise. The per-antenna SNR is \(\text{SNR} = P/(N_t \sigma^2)\).

The fundamental trade-off in MIMO systems is between spatial multiplexing gain (transmitting multiple independent streams) and diversity gain (transmitting redundant copies to combat fading). Zheng and Tse \cite{zheng2003diversity} formalized this as the diversity-multiplexing tradeoff (DMT): for an \(N_t \times N_r\) system at multiplexing gain \(r\), the maximum achievable diversity order is \(d(r) = (N_t - r)(N_r - r)\). In the \(2 \times 2\) case at full spatial multiplexing (\(r = 2\)), \(d = 0\) --- meaning no diversity protection and the system is interference-limited. This makes the choice of equalization method critical.

\subsection{Equalization Methods}\label{sec:equalization-methods}

Equalization at the receiver aims to recover \(\mathbf{x}\) from \(\mathbf{y}\) in the presence of inter-stream interference. Three methods of increasing sophistication are employed in this thesis:

\textbf{Zero-Forcing (ZF).} The ZF equalizer applies the pseudo-inverse of the channel:

\begin{equation}\protect\phantomsection\label{eq:zf-equalizer}{\hat{\mathbf{x}}_{\text{ZF}} = (\mathbf{H}^H\mathbf{H})^{-1}\mathbf{H}^H\mathbf{y} = \mathbf{x} + (\mathbf{H}^H\mathbf{H})^{-1}\mathbf{H}^H\mathbf{n}
}\end{equation} \par{\small\textit{Source: \cite{tse2005fundamentals, tse1999linear}}}

This completely eliminates inter-stream interference but amplifies noise. The post-equalization SNR for stream \(k\) is:

\begin{equation}\protect\phantomsection\label{eq:zf-snr}{\text{SNR}_k^{\text{ZF}} = \frac{\text{SNR}}{[(\mathbf{H}^H\mathbf{H})^{-1}]_{kk}}
}\end{equation} \par{\small\textit{Source: \cite{tse2005fundamentals, tse1999linear}}}

When \(\mathbf{H}\) is ill-conditioned (i.e., its singular values are disparate), the diagonal elements of \((\mathbf{H}^H\mathbf{H})^{-1}\) become large, severely degrading performance. ZF achieves a diversity order of \(d = N_r - N_t + 1 = 1\) for the \(2 \times 2\) case \cite{tse2005fundamentals}, meaning it provides minimal diversity protection and its BER decays as \(1/\text{SNR}\).

\textbf{Minimum Mean Square Error (MMSE).} The MMSE equalizer minimizes \(\mathbb{E}[\|\hat{\mathbf{x}} - \mathbf{x}\|^2]\), yielding the Wiener filter:

\begin{equation}\protect\phantomsection\label{eq:mmse-equalizer}{\hat{\mathbf{x}}_{\text{MMSE}} = (\mathbf{H}^H\mathbf{H} + \sigma^2\mathbf{I})^{-1}\mathbf{H}^H\mathbf{y}
}\end{equation} \par{\small\textit{Source: \cite{tse2005fundamentals, tse1999linear}}}

The regularization term \(\sigma^2\mathbf{I}\) prevents noise amplification by biasing the estimate toward zero when the channel is weak. The post-equalization SINR for stream \(k\) is:

\begin{equation}\protect\phantomsection\label{eq:mmse-sinr}{\text{SINR}_k^{\text{MMSE}} = \frac{1}{[(\mathbf{H}^H\mathbf{H} + \sigma^2\mathbf{I})^{-1}]_{kk}} - 1
}\end{equation} \par{\small\textit{Source: \cite{tse2005fundamentals, tse1999linear}}}

At low SNR, MMSE significantly outperforms ZF because it avoids noise amplification; at high SNR (\(\sigma^2 \to 0\)), the MMSE filter converges to the ZF solution. Crucially, MMSE achieves the same diversity order as ZF (\(d = 1\)) but with a superior coding gain, meaning it provides a constant SNR advantage across the entire operating range \cite{tse1999linear}.

\textbf{Successive Interference Cancellation (SIC).} SIC is a non-linear detection technique that exploits the layered structure of spatial multiplexing. The MMSE-ordered V-BLAST algorithm \cite{wolniansky1998vblast}, \cite{foschini1996layered} proceeds as follows:

\begin{enumerate}
\def\labelenumi{\arabic{enumi}.}
\tightlist
\item
  \textbf{Order} streams by post-MMSE SINR: detect the strongest stream first.
\item
  \textbf{Detect} stream \(k\) using the MMSE filter on the residual observation.
\item
  \textbf{Cancel} the contribution of stream \(k\): \(\mathbf{y} \leftarrow \mathbf{y} - \mathbf{h}_k \hat{x}_k\).
\item
  \textbf{Repeat} for the remaining stream(s) with a reduced-dimension channel.
\end{enumerate}

The key advantage is that after perfect cancellation, the second stream sees an interference-free channel with MMSE filtering, yielding higher post-detection SINR. For a \(2 \times 2\) system, the first detected stream achieves diversity order \(d_1 = N_r - N_t + 1 = 1\), while the second stream (after cancellation) achieves \(d_2 = N_r = 2\) \cite{zheng2003diversity}. This gives an average diversity order strictly better than linear equalizers.

However, SIC is vulnerable to \textbf{error propagation}: if the first stream is decoded incorrectly, the residual interference from the incorrect cancellation degrades the second stream. This effect is most pronounced at low SNR, where the first-stream BER is high. MMSE-ordered detection mitigates this by choosing the most reliable stream first, but cannot eliminate it entirely \cite{loyka2004performance}.

\subsection{MIMO with Relay Processing}\label{sec:mimo-with-relay-processing}

In the two-hop relay architecture studied in this thesis, MIMO equalization operates at the destination \emph{after} Hop 2, while the neural network relay operates at the intermediate node \emph{after} Hop 1. These are independent signal-processing stages that solve different problems: the relay denoises the scalar (SISO) or per-antenna signal, while the equalizer separates spatially multiplexed streams. The combined system applies relay processing first (per-antenna denoising), then MIMO equalization at the destination.

Combining AI-based relay processing with MIMO equalization has not been studied in the literature. A key question is whether AI relays that excel in the SISO setting maintain their advantage when cascaded with linear or non-linear MIMO equalizers. This thesis evaluates all nine relay strategies across all three equalization methods, providing the first systematic study of this interaction.

\section{Research Gap and Motivation}\label{sec:research-gap-and-motivation}

Despite growing interest in AI for wireless communication, a systematic comparison of relay processing paradigms --- spanning classical signal processing, supervised learning, generative modeling, and modern sequence architectures --- is absent from the literature. The following specific gaps motivate this thesis:

\textbf{Gap 1: No cross-paradigm relay comparison.} Existing work on AI-based relay processing focuses on individual architectures in isolation. Ye et al.~\cite{ye2018power} demonstrated DNNs for OFDM detection, Dorner et al.~\cite{dorner2018deep} explored autoencoders for end-to-end communication, and Samuel et al.~\cite{samuel2019learning} applied unfolded networks to MIMO detection. However, no study systematically compares supervised feedforward networks, variational autoencoders, adversarial networks, attention-based Transformers, and state space models (Mamba) for the same relay denoising task under controlled conditions. Without such a comparison, practitioners cannot make informed architecture selections.

\textbf{Gap 2: Model complexity--performance relationship is unknown.} It is commonly assumed that larger neural networks yield better performance. However, the relay denoising task --- recovering a BPSK symbol from a noisy observation --- is a low-dimensional mapping problem. The theoretical minimum description length for this mapping is small (a soft threshold function), suggesting that large models may overfit. No prior work has rigorously characterized the complexity--performance trade-off for relay processing, nor has a normalized (equal-parameter) comparison been conducted to isolate architectural effects from capacity effects.

\textbf{Gap 3: Limited channel diversity in evaluations.} Most AI-for-wireless studies evaluate on a single channel model (typically AWGN or Rayleigh). The performance of AI relays under Rician fading, MIMO spatial multiplexing, and different equalization strategies has not been studied. Whether AI relays trained on one channel generalize to others --- and whether their advantage over classical methods persists across diverse propagation conditions --- are open questions.

\textbf{Gap 4: State space models for physical-layer signal processing.} The Mamba architecture \cite{gu2024mamba} and its Mamba-2 SSD variant \cite{dao2024transformers} have demonstrated strong performance in NLP and genomics, but their application to physical-layer wireless communication is entirely unexplored. The theoretical connection between SSMs and classical adaptive filters makes this a natural research direction, but empirical validation is needed.

\textbf{Gap 5: AI relay--MIMO equalization interaction.} While MIMO equalization (ZF, MMSE, SIC) and AI-based processing have been studied independently, their \textbf{joint} effect in a relay pipeline --- where the relay denoises the signal before equalization separates the spatial streams --- has not been investigated. Whether the gains from better relay processing and better equalization are additive, synergistic, or partially redundant is unknown.

The following table summarizes the positioning of this work relative to the literature:

\begin{longtable}[]{@{}
  >{\raggedright\arraybackslash}p{(\linewidth - 4\tabcolsep) * \real{0.3333}}
  >{\raggedright\arraybackslash}p{(\linewidth - 4\tabcolsep) * \real{0.3333}}
  >{\raggedright\arraybackslash}p{(\linewidth - 4\tabcolsep) * \real{0.3333}}@{}}
\caption{Research positioning: comparison of this thesis against prior work in deep-learning-based relay communication.}\label{tbl:table17}\tabularnewline
\toprule\noalign{}
\begin{minipage}[b]{\linewidth}\raggedright
Aspect
\end{minipage} & \begin{minipage}[b]{\linewidth}\raggedright
Prior Work
\end{minipage} & \begin{minipage}[b]{\linewidth}\raggedright
This Thesis
\end{minipage} \\
\midrule\noalign{}
\endfirsthead
\toprule\noalign{}
\begin{minipage}[b]{\linewidth}\raggedright
Aspect
\end{minipage} & \begin{minipage}[b]{\linewidth}\raggedright
Prior Work
\end{minipage} & \begin{minipage}[b]{\linewidth}\raggedright
This Thesis
\end{minipage} \\
\midrule\noalign{}
\endhead
\bottomrule\noalign{}
\endlastfoot
Relay strategies compared & 1--2 (typically AF vs.~DF, or one AI method) & 9 (2 classical + 7 AI) \\
AI paradigms & Single (e.g., supervised only) & 4 (supervised, generative, adversarial, sequential) \\
Channel models & 1--2 (AWGN, sometimes Rayleigh) & 6 (AWGN + Rayleigh + Rician SISO; MIMO ZF/MMSE/SIC) \\
Complexity analysis & None & Normalized 3K comparison + inverted-U study \\
State space models & None for relay/physical-layer & Mamba S6, Mamba-2 SSD, with crossover benchmark \\
Statistical rigor & Often missing & Wilcoxon test, 95\% CI, 10 trials per point \\
\end{longtable}

This thesis addresses all five gaps through a unified framework that implements and compares nine relay strategies across six channel/topology configurations, with both original and normalized parameter counts, full statistical analysis, and a dedicated complexity study.

\begin{center}\rule{0.5\linewidth}{0.5pt}\end{center}

\section{Theoretical Foundations: Channel Models}\label{sec:theoretical-foundations-channel-models}

This section presents the theoretical BER analysis for each channel model used in this study, derives the closed-form expressions, and validates them against Monte Carlo simulation. The theoretical--simulative comparison serves two purposes: (i) it verifies the correctness of the simulation framework, and (ii) it establishes the baseline performance that AI relays must beat.

\subsection{AWGN Channel --- Theoretical Analysis}\label{sec:awgn-channel-theoretical-analysis}

The AWGN channel adds zero-mean Gaussian noise to the transmitted signal:

\begin{equation}\protect\phantomsection\label{eq:awgn-channel}{y = x + n, \quad n \sim \mathcal{N}(0, \sigma^2), \quad \sigma^2 = \frac{P_s}{\text{SNR}_{\text{linear}}}
}\end{equation} \par{\small\textit{Source: \cite{proakis2008digital}}}

where \(P_s = \mathbb{E}[|x|^2] = 1\) for unit-power BPSK symbols. The noise variance is inversely proportional to the linear SNR.

\textbf{Single-hop BER.} For BPSK \(\pm 1\) symbols and real-valued AWGN with variance \(\sigma^2 = 1/\text{SNR}\), an error occurs when the noise pushes the received sample past the decision boundary at the origin. The theoretical BER is:

\begin{equation}\protect\phantomsection\label{eq:awgn-ber}{P_e^{\text{AWGN}} = Q\!\left(\frac{1}{\sigma}\right) = Q\left(\sqrt{\text{SNR}}\right) = \frac{1}{2}\,\text{erfc}\!\left(\sqrt{\frac{\text{SNR}}{2}}\right)
}\end{equation} \par{\small\textit{Source: \cite{proakis2008digital}}}

where \(Q(x) = \frac{1}{2}\text{erfc}(x/\sqrt{2})\) is the Gaussian Q-function and \(\text{SNR} = P_s / \sigma^2\) is the ratio of signal power to noise variance.

\begin{quote}
\textbf{Noise-convention note.} The standard textbook result \(P_e = Q(\sqrt{2 E_b/N_0})\) assumes that the one-sided noise PSD is \(N_0\), giving a baseband noise variance of \(N_0/2\) per real dimension. In our AWGN implementation the noise is purely real with variance \(\sigma^2 = P_s/\text{SNR}\), so \(N_0 = 2\sigma^2 = 2P_s/\text{SNR}\) and \(E_b/N_0 = \text{SNR}/2\). Substituting yields \(Q(\sqrt{2 \cdot \text{SNR}/2}) = Q(\sqrt{\text{SNR}})\), which is the expression used throughout this work.

By contrast, the fading channels add \textbf{complex} Gaussian noise with total power \(1/\text{SNR}\) (i.e.~\(\sigma_I^2 = \sigma_Q^2 = 1/(2\,\text{SNR})\)) and extract the real part after ZF equalization, halving the effective noise variance per decision dimension. This naturally introduces a factor of 2 inside the Q-function argument for all fading-channel BER expressions (Sections 6.2.2--6.2.5), making them consistent with the standard textbook forms.
\end{quote}

\textbf{Two-hop DF BER.} For a decode-and-forward relay with equal-SNR hops, an error occurs at the destination if exactly one hop introduces an error:

\begin{equation}\protect\phantomsection\label{eq:df-ber-awgn}{P_e^{\text{DF}} = P_{e,1} + P_{e,2} - 2P_{e,1}P_{e,2}
}\end{equation} \par{\small\textit{Source: \cite{laneman2004cooperative}}}

With equal hops (\(P_{e,1} = P_{e,2} = P_e^{\text{AWGN}}\)), this becomes \(P_e^{\text{DF}} = 2P_e(1 - P_e)\). At high SNR, \(P_e \ll 1\) and \(P_e^{\text{DF}} \approx 2P_e\), i.e., the two-hop penalty is approximately a factor of 2 in BER.

\textbf{Two-hop AF BER.} For amplify-and-forward with equal-SNR hops, the effective end-to-end SNR is:

\begin{equation}\protect\phantomsection\label{eq:af-snr-eff}{\text{SNR}_{\text{eff}}^{\text{AF}} = \frac{\text{SNR}_1 \cdot \text{SNR}_2}{\text{SNR}_1 + \text{SNR}_2 + 1} = \frac{\gamma^2}{2\gamma + 1}
}\end{equation} \par{\small\textit{Source: \cite{laneman2004cooperative}}}

where \(\gamma = \text{SNR}\) is the per-hop SNR. The resulting BER is \(P_e^{\text{AF}} = Q\left(\sqrt{\text{SNR}_{\text{eff}}}\right)\) (same real-noise convention as the single-hop case). At high SNR, \(\text{SNR}_{\text{eff}} \approx \gamma/2\), confirming the well-known 3 dB penalty of AF relaying.


\subsection{Rayleigh Fading Channel --- Theoretical Analysis}\label{sec:rayleigh-fading-channel-theoretical-analysis}

The Rayleigh fading channel models non-line-of-sight (NLOS) propagation where the signal undergoes multiplicative fading:

\begin{equation}\protect\phantomsection\label{eq:rayleigh-channel}{y = hx + n, \quad h \sim \mathcal{CN}(0, 1), \quad n \sim \mathcal{CN}(0, \sigma^2)
}\end{equation} \par{\small\textit{Source: \cite{proakis2008digital}}}

The fading coefficient \(h\) is a circularly-symmetric complex Gaussian random variable, so its magnitude \(|h|\) follows a Rayleigh distribution:

\begin{equation}\protect\phantomsection\label{eq:rayleigh-pdf}{f_{|h|}(r) = 2r \cdot e^{-r^2}, \quad r \geq 0
}\end{equation} \par{\small\textit{Source: \cite{proakis2008digital}}}

with \(\mathbb{E}[|h|^2] = 1\) (unit average power). The instantaneous SNR after equalization (\(\hat{x} = y/h\)) becomes \(\gamma_h = |h|^2 \cdot \text{SNR}\), which is exponentially distributed.

\textbf{Single-hop BER.} Averaging the conditional BER \(P_e(\gamma_h) = Q(\sqrt{2\gamma_h})\) over the exponential distribution of \(\gamma_h\) yields the closed-form {[}21, Eq. 14-4-15{]}:

\begin{equation}\protect\phantomsection\label{eq:rayleigh-ber}{P_e^{\text{Rayleigh}} = \frac{1}{2}\left(1 - \sqrt{\frac{\bar{\gamma}}{1 + \bar{\gamma}}}\right)
}\end{equation} \par{\small\textit{Source: \cite{proakis2008digital}}}

where \(\bar{\gamma} = \text{SNR}\) is the average SNR. At high SNR (\(\bar{\gamma} \gg 1\)):

\begin{equation}\protect\phantomsection\label{eq:rayleigh-ber-approx}{P_e^{\text{Rayleigh}} \approx \frac{1}{4\bar{\gamma}}
}\end{equation} \par{\small\textit{Source: \cite{proakis2008digital}}}

This \(1/\text{SNR}\) decay is fundamentally slower than the exponential decay of AWGN (\(Q(\sqrt{\gamma}) \sim e^{-\gamma/2}\)), explaining why Rayleigh fading is significantly more challenging. The channel is \textbf{diversity-limited}: deep fades (where \(|h| \approx 0\)) cause errors regardless of the average SNR. This is a primary motivation for MIMO systems, which exploit spatial diversity to combat fading.

\textbf{Two-hop DF BER.} The two-hop DF relay over Rayleigh fading follows the same composition rule as AWGN:

\begin{equation}\protect\phantomsection\label{eq:df-ber-rayleigh}{P_e^{\text{DF,Rayleigh}} = 2P_e^{\text{Rayleigh}}(1 - P_e^{\text{Rayleigh}})
}\end{equation} \par{\small\textit{Source: \cite{laneman2004cooperative}}}


\subsection{Rician Fading Channel --- Theoretical Analysis}\label{sec:rician-fading-channel-theoretical-analysis}

The Rician fading channel models environments with a dominant line-of-sight (LOS) component alongside scattered multipath:

\begin{equation}\protect\phantomsection\label{eq:rician-channel}{h = \sqrt{\frac{K}{K+1}} e^{j\theta} + \sqrt{\frac{1}{K+1}} h_{\text{scatter}}, \quad h_{\text{scatter}} \sim \mathcal{CN}(0, 1)
}\end{equation} \par{\small\textit{Source: \cite{simon2005digital}}}

The K-factor is the ratio of LOS power to scatter power. The fading amplitude \(|h|\) follows a Rician distribution:

\begin{equation}\protect\phantomsection\label{eq:rician-pdf}{f_{|h|}(r) = \frac{r}{\sigma^2} \exp\left(-\frac{r^2 + \nu^2}{2\sigma^2}\right) I_0\left(\frac{r\nu}{\sigma^2}\right)
}\end{equation} \par{\small\textit{Source: \cite{simon2005digital}}}

where \(\nu = \sqrt{K/(K+1)}\) is the LOS amplitude, \(\sigma^2 = 1/(2(K+1))\) is the scatter variance per component, and \(I_0(\cdot)\) is the modified Bessel function of the first kind, order zero.

\textbf{Special cases.} When \(K = 0\), the LOS component vanishes and the Rician channel degenerates to Rayleigh. As \(K \to \infty\), the channel approaches AWGN (no fading). Thus the Rician model interpolates between the two extreme cases, with \(K=3\) (used in this study) representing a moderate LOS environment.

\textbf{Single-hop BER.} The average BER for BPSK over a Rician channel is obtained via the moment-generating function (MGF) approach \cite{simon2005digital}:

\begin{equation}\protect\phantomsection\label{eq:rician-ber-mgf}{P_e^{\text{Rician}} = \frac{1}{\pi} \int_0^{\pi/2} M_{\gamma}\left(\frac{-1}{\sin^2\theta}\right) d\theta
}\end{equation} \par{\small\textit{Source: \cite{simon2005digital}}}

where the MGF of the instantaneous SNR \(\gamma\) under Rician fading is:

\begin{equation}\protect\phantomsection\label{eq:rician-mgf}{M_{\gamma}(s) = \frac{1+K}{1+K - s\bar{\gamma}} \cdot \exp\left(\frac{Ks\bar{\gamma}}{1+K-s\bar{\gamma}}\right)
}\end{equation} \par{\small\textit{Source: \cite{simon2005digital}}}

This integral is evaluated numerically. The resulting BER falls between the AWGN and Rayleigh curves, with the position determined by the K-factor.

\textbf{Two-hop DF BER.} As with other channels, \(P_e^{\text{DF,Rician}} = 2P_e^{\text{Rician}}(1 - P_e^{\text{Rician}})\).


\subsection{Fading Coefficient Distributions}\label{sec:fading-coefficient-distributions}

The statistical behavior of the fading coefficient \(|h|\) determines the severity of fading for each channel model. Figure~\ref{fig:fig4} shows the probability density function (PDF) and cumulative distribution function (CDF) of \(|h|\) for Rayleigh and Rician fading with various K-factors.

\textbf{Key observations from the fading PDFs:}

\begin{itemize}
\tightlist
\item
  \textbf{Rayleigh (\(K=0\)):} The PDF is maximized at \(|h| = 1/\sqrt{2} \approx 0.707\) and has a heavy tail toward zero, meaning deep fades are common. The probability of a deep fade (\(|h| < 0.3\)) is approximately 8.6\%.
\item
  \textbf{Rician \(K=1\):} The LOS component shifts the PDF peak toward higher amplitudes, reducing the probability of deep fades.
\item
  \textbf{Rician \(K=3\):} The distribution becomes more concentrated around the LOS amplitude (\(\nu \approx 0.87\)). Deep fade probability drops to approximately 0.3\%.
\item
  \textbf{Rician \(K=10\):} Approaches a near-deterministic channel; the PDF becomes sharply peaked around \(|h| \approx 0.95\). Fading is negligible.
\end{itemize}

The CDF plot directly shows the \textbf{outage probability} \(P(|h| \leq x)\): for a given threshold \(x\), the CDF value gives the probability that the fading amplitude falls below that threshold. Rayleigh has the highest outage probability at any threshold, confirming its role as the worst-case fading model.


\subsection{2×2 MIMO Channel --- Theoretical Analysis}\label{sec:mimo-channel-theoretical-analysis}

The 2×2 MIMO spatial multiplexing system transmits two independent BPSK streams simultaneously:

\begin{equation}\protect\phantomsection\label{eq:mimo-received-2x2}{\mathbf{y} = \mathbf{H}\mathbf{x} + \mathbf{n}, \quad \mathbf{H} \in \mathbb{C}^{2 \times 2}, \quad H_{ij} \sim \mathcal{CN}(0, 1), \quad \mathbf{n} \sim \mathcal{CN}(\mathbf{0}, \sigma^2\mathbf{I})
}\end{equation} \par{\small\textit{Source: \cite{tse2005fundamentals}}}

The theoretical per-stream BER depends on the equalization technique:

\textbf{ZF equalization.} After ZF equalization (\(\hat{\mathbf{x}} = \mathbf{H}^{-1}\mathbf{y}\)), the effective noise on stream \(k\) has variance \(\sigma^2 [\mathbf{H}^{-1}(\mathbf{H}^{-1})^H]_{kk}\). For a 2×2 system with i.i.d. Rayleigh fading, each post-ZF stream sees an effective diversity order of \(n_R - n_T + 1 = 1\), identical to SISO Rayleigh. Therefore:

\begin{equation}\protect\phantomsection\label{eq:zf-ber-mimo}{P_e^{\text{ZF}} \approx \frac{1}{2}\left(1 - \sqrt{\frac{\bar{\gamma}}{1 + \bar{\gamma}}}\right)
}\end{equation} \par{\small\textit{Source: \cite{tse2005fundamentals, zheng2003diversity}}}

This is the same expression as SISO Rayleigh --- ZF spatial multiplexing provides no diversity gain for a square (\(n_T = n_R\)) system.

\textbf{MMSE equalization.} The MMSE filter \(\mathbf{W} = (\mathbf{H}^H\mathbf{H} + \sigma^2\mathbf{I})^{-1}\mathbf{H}^H\) provides a noise-regularized estimate. The post-MMSE SINR exceeds the post-ZF SNR because the regularization prevents extreme noise amplification when \(\mathbf{H}\) is ill-conditioned. The exact BER analysis requires integration over the joint distribution of post-MMSE SINRs, which does not admit a simple closed form for 2×2 systems. An effective SNR approximation yields:

\begin{equation}\protect\phantomsection\label{eq:mmse-ber-mimo}{P_e^{\text{MMSE}} \approx \frac{1}{2}\left(1 - \sqrt{\frac{\gamma_{\text{eff}}}{1 + \gamma_{\text{eff}}}}\right), \quad \gamma_{\text{eff}} \approx \bar{\gamma} \cdot \left(1 + \frac{1}{\bar{\gamma} + 1}\right)
}\end{equation} \par{\small\textit{Source: \cite{tse2005fundamentals, tse1999linear}}}

The MMSE gain over ZF is most significant at low SNR (where the regularization term dominates) and diminishes at high SNR (where \(\sigma^2 \to 0\) and MMSE converges to ZF).

\textbf{SIC equalization.} The MMSE-SIC (V-BLAST) receiver detects the stronger stream first via MMSE, makes a hard decision, cancels its contribution, and detects the remaining stream interference-free. When the first decision is correct (which is the common case, since the stronger stream has higher SINR), the second stream sees no inter-stream interference, effectively achieving the single-stream MMSE bound. The improvement over linear MMSE comes from eliminating the interference term for the second stream.


\subsection{Channel Model Validation}
The analytical error probability models described above have been rigorously validated against the Monte Carlo simulation framework developed for this thesis. Detailed comparisons between the theoretical bounds and empirical bit error rates for single-hop transmission, along with their comprehensive discussion, are deferred to the experimental results in Chapter~\ref{sec:experiments} (Section~\ref{sec:channel-model-validation}).

